\documentclass[11pt,a4paper]{article}
\usepackage[utf8]{inputenc}
\usepackage[T1]{fontenc}
\usepackage[margin=1in]{geometry}
\usepackage{parskip}
\usepackage{hyperref}
\usepackage{xurl}
\hypersetup{hidelinks}
\pagestyle{empty}
\begin{document}

\textbf{To:} The Editors, \emph{Quantum Machine Intelligence}\\
\textbf{Re:} Submission of original research article\\
\textbf{Topic:} Quantum Computing for AI

\bigskip
Dear Editors,

We are pleased to submit our manuscript, \textbf{``How Quantum Is the Advantage? A Fair,
Calibration- and Noise-Aware Benchmark and Attribution Audit of Quantum Machine Learning for
Network Intrusion Detection,''} for consideration as an original research article in
\emph{Quantum Machine Intelligence}.

\textbf{The question.} Applied quantum machine learning reports near-perfect accuracies across
many application domains, network intrusion detection among them, with figures in the 97 to
99.9\% range. The field's more careful studies have begun to question those numbers, showing that
well-tuned classical models remain competitive on tabular data and that apparent quantum gains may
be artefacts of the classical dimensionality reduction and implicit regularisation that every
hybrid pipeline contains. What the literature lacks is not another accuracy record but a
\emph{method} for deciding, for a given result, how much of the measured advantage is actually
attributable to the quantum component. That is what this paper supplies.

\textbf{What we contribute.} The methodological core is a \textbf{quantum-attribution audit}: for
every quantum model we construct a parameter-matched classical control sharing the identical
classical front-end; for every quantum kernel we construct a random-feature kernel as its explicit
classical surrogate; and we run a regularisation sweep to test whether comparable regularisation
closes any observed gap. Differences are reported per variant and metric with significance testing
and Benjamini-Hochberg control of the false-discovery rate across the whole comparison family. The
audit is general: it applies to any claim of quantum advantage in a hybrid learning pipeline, and
we demonstrate it on a benchmark of two quantum model families (hybrid variational circuits and
quantum-kernel SVMs) against five honestly-tuned classical baselines, across four standard
datasets, under one leakage-controlled protocol with an equal-budget feature view, imbalance- and
calibration-aware metrics, and a simulated NISQ noise sweep.

\textbf{What we find.} Tuned classical models match or exceed the quantum models on aggregate
detection on every dataset, and the audit attributes this to classical preprocessing and
regularisation rather than to quantum effects. Two advantages survive FDR control. The
quantum-kernel SVM significantly out-ranks its direct classical surrogate on AUPRC ($q=0.011$) and
ROC-AUC ($q=0.018$); because the random-feature kernel is the explicit classical analogue of the
quantum feature map, this is the sharpest isolation of a genuinely quantum contribution in the
study, and it is the one result that also survives the more conservative family-wise Holm
correction. Separately, a small four-qubit hybrid out-detects the best classical baseline at the
1\% false-positive operating point on a distribution-shifted task ($p=0.005$, $q=0.030$); we report
plainly that this one does not survive Holm.

\textbf{Fit for the journal.} The work sits in the journal's Quantum Computing for AI perspective
and speaks directly to its stated emphasis on benchmarking and empirical evaluation. It is
addressed to readers who want to know whether reported quantum advantages hold up under matched
classical controls, which is a question for this readership rather than for an applications venue.
We also believe the audit is reusable well beyond intrusion detection.

\textbf{Limitations we state up front.} All quantum results are simulated; we have no QPU access
and treat the simulated numbers as an optimistic bound. Our seed budget is small and uneven (five
seeds for the classical baselines on one dataset, three elsewhere and for all quantum runs, two for
the noise sweep), which we report explicitly and which is why the operating-point claims rest on a
paired bootstrap rather than on seed-level tests. The attribution audit is decomposed on the most
discriminative dataset. We prefer to surface these limits than have them inferred.

\textbf{Reproducibility.} The code, configurations, fixed seeds, the scripts that reconstruct our
exact splits, and the per-run result records are released publicly; every reported number is
regenerable from the released harness.

This manuscript is original, has not been published previously, and is not under consideration
elsewhere. All authors have approved the submission and declare no competing interests.

Thank you for your consideration.

\bigskip
Sincerely,

\smallskip
Syeda Anshrah Gillani (corresponding author), Hamdard University, Karachi, Pakistan\\
\texttt{SyedaAnshrah16@gmail.com}\\[2pt]
Mirza Samad Ahmed Baig (corresponding author), Fandaqah, Al Khobar, Saudi Arabia\\
\texttt{MirzaSamadcontact@gmail.com}\\[2pt]
Shahid Munir Shah, SZABIST University, Pakistan\\[2pt]
Asher Ali, Fandaqah, Al Khobar, Saudi Arabia\\[2pt]
Hamzah Siddiqui, Fandaqah, Al Khobar, Saudi Arabia and Hamdard University, Karachi, Pakistan

\end{document}
